\abstracten

Metadata operations---file creation, deletion, renaming, and the like---account for 60\%--80\% of file system I/O, and the storage structure and operational efficiency of metadata have a decisive impact on overall system throughput. Traditional file systems manage directory entries and inode indexes with B+ trees, where every insertion or deletion requires in-place page modification and journal synchronization, causing heavy random I/O and directory-level lock contention under high-concurrency small-file workloads. The fundamental reason is that B+ trees operate at page granularity: even a single-entry change forces a full page read and rewrite. Key-value (KV) storage refines this granularity to the individual record: each directory entry becomes an independent key-value pair, and an insertion or deletion reduces to a single append write with no page-structure maintenance involved. LSM-tree (Log-Structured Merge-tree) based KV engines build on this by aggregating many append writes in memory and flushing them to disk sequentially in batches, reducing write cost by roughly an order of magnitude.

Based on these observations, this work designs and implements RucksFS---a user-space file system backed by RocksDB and exposed as a POSIX mount via FUSE---with a complete set of file operations including \texttt{create}, \texttt{mkdir}, \texttt{unlink}, \texttt{rmdir}, \texttt{rename}, \texttt{readdir}, and \texttt{setattr}. Implementing these operations on KV storage requires answering three design questions: how to model metadata as key-value structures, how to control parent-directory write amplification, and how to guarantee correctness under concurrent operations.

For storage modeling, the system adopts a multi-column-family architecture, placing inode attributes and directory entries in separate column families with individually tuned Bloom filters and compaction policies. The key design for the directory-entry column family follows an edge-centric principle: the parent inode number and the child filename are concatenated in big-endian order as the key, so every directory entry maps directly to a parent--child edge in the directory tree. File operations accordingly become edge manipulations---\texttt{create} inserts an edge, \texttt{unlink} removes one, \texttt{rename} deletes the old edge and inserts a new one---all in constant time. For comparison, full-path keying schemes must recursively rewrite every key in the subtree on a \texttt{rename}, at a cost that grows linearly with subtree size. In addition, all child edges of the same parent are naturally adjacent in the key space, so directory traversal and single-entry lookup can both be completed efficiently via prefix scan and point query.

For write amplification, POSIX semantics require that \texttt{create} and \texttt{unlink} update the parent directory's timestamps and link count, making a busy parent a write hot spot under high concurrency; read--modify--write on such a hot spot would severely throttle throughput. A delta update and lazy folding mechanism addresses this: on the write path, each operation only appends a small delta record without touching the base inode value, turning in-place rewrites into append writes. On the read path, the base value and all outstanding deltas are merged on the fly before being returned, ensuring that callers always observe up-to-date attributes. A background thread periodically scans accumulated deltas and folds them back into the base value once a threshold is reached, keeping delta counts bounded to prevent read-path degradation.

For concurrency safety, when multiple threads operate on the same directory, an inherent race window exists between checking and modifying---\texttt{create} must first confirm no duplicate exists before inserting, \texttt{rmdir} must first confirm the directory is empty before deleting, and \texttt{rename} additionally involves multi-key atomic modifications across directories. RocksDB Pessimistic Concurrency Control (PCC) transactions lock the involved keys one by one, wrapping both checks and mutations inside a single atomic operation to eliminate these races.

For correctness, a full pjdfstest run in a distributed deployment (client and MetadataServer on separate nodes) executes 398 cases. Partitioned by the scope of this work: 167 cases directly or transitively depend on \texttt{mknod}/\texttt{mkfifo}, which this work deliberately leaves unimplemented; 1 case exceeds the DataServer's current 64\,MiB per-file development limit; 48 cases are skipped by the pjdfstest harness itself (\texttt{utimensat}, \texttt{remount}, etc., unrelated to the metadata design). The remaining 182 cases, which constitute the in-scope set for this work, all pass---a 100\% pass rate within scope. For performance, mdtest benchmarks run on a cluster of one 64-core server and up to 96 two-core clients. For $N \in \{2, 8, 32\}$, RucksFS is compared head-to-head with NFS v4.2 and JuiceFS+TiKV: on \texttt{create}, \texttt{remove}, and \texttt{stat}, RucksFS is about 1.5--1.7$\times$, 1.8--2.2$\times$, and 1.2$\times$ faster than JuiceFS+TiKV respectively---partly because a single-node direct-to-RocksDB design avoids TiKV's Raft and two-phase commit overhead, partly because the DeltaOp mechanism reduces contention on shared parent directories. This places the RucksFS metadata implementation on par with mainstream KV-backed file systems. Against NFS v4.2, with all client-side attribute caches disabled for a fair MDS-path comparison, RucksFS is actually slower at $N=2$ (each \texttt{create} now requires three RPCs instead of one), but at $N=32$, once \texttt{i\_rwsem} contention kicks in on the ext4 side, RucksFS overtakes NFS by about 4.6$\times$ on \texttt{create}, illustrating that the KV approach scales better than B-tree plus directory lock under shared hot directories. At $N \geq 64$ neither baseline sustains reliable numbers, so the high-concurrency comparison is done internally between RucksFS's Delta and NoDelta builds, sweeping $\mathrm{np} = N \times r$ in two dimensions: once $\mathrm{np}$ crosses 96, NoDelta enters a transaction-conflict loop and stabilises around 12{,}000\,ops/s, while Delta continues to scale to 38{,}000\,ops/s, yielding a create ratio of up to 2.96$\times$ and a remove ratio of up to 2.63$\times$. Under independent-directory mode the two variants are indistinguishable, confirming that the gain originates from eliminating write conflicts on shared parent directories.

\vspace{5pt}

\keywordsen File System, Metadata Management, Key-Value Storage, LSM-tree, Delta Update, Pessimistic Concurrency Control, FUSE
